\documentclass[../main.tex]{subfiles}
\begin{document}

\begin{center}
    \Large{\textbf{ABSTRACT}}\\
\end{center}
\vspace{1cm}
Fashion sketch-to-image generation has become an important research topic due to its potential to assist designers in rapidly transforming conceptual sketches into realistic garment visualizations. Such systems can significantly accelerate the fashion design workflow by reducing the time required for manual rendering and enabling fast design iteration. Existing approaches to this problem mainly rely on image-to-image translation techniques, particularly Generative Adversarial Networks (GANs), which learn the mapping between sketch and photo domains. Recent studies have demonstrated impressive visual quality; however, most state-of-the-art models are computationally expensive, contain a large number of parameters, and require substantial inference time. These limitations restrict their applicability in interactive design environments where near real-time response and efficient deployment are essential. Furthermore, publicly available fashion sketch datasets are often limited in scale, lack attribute annotations, or are not specifically designed to support lightweight generation models.

To address these challenges, this thesis focuses on developing a lightweight GAN-based framework for fashion sketch-to-image translation that prioritizes efficiency while maintaining stable image quality. This direction was selected because GANs offer a favorable balance between visual fidelity and computational cost compared with many recent large-scale generative models, making them more suitable for real-time design assistance applications. The proposed solution consists of two main components. First, a large-scale paired dataset is constructed by leveraging publicly available fashion image datasets, particularly Fashionpedia, and applying edge-based sketch generation techniques to create corresponding sketch representations. The resulting dataset is further enriched with fashion attribute annotations to facilitate future research and model development. Second, a lightweight GAN architecture is designed to generate realistic fashion images from sketch inputs while reducing model complexity and inference latency.

The main contributions of this thesis are the construction of a new dataset containing approximately 40,000 paired sketch–fashion image samples with annotated attributes and the development of a lightweight GAN-based model capable of generating realistic fashion images from sketches. The final system demonstrates the feasibility of supporting rapid design visualization and lays the foundation for future interactive fashion design applications that require fast generation and iterative editing capabilities.
% Sinh viên viết tóm tắt ĐATN của mình trong mục này, với 200 đến 350 từ. Theo trình tự, các nội dung tóm tắt cần có: (i) Giới thiệu vấn đề (tại sao có vấn đề đó, hiện tại được giải quyết chưa, có những hướng tiếp cận nào, các hướng này giải quyết như thế nào, hạn chế là gì), (ii) Hướng tiếp cận sinh viên lựa chọn là gì, vì sao chọn hướng đó, (iii) Tổng quan giải pháp của sinh viên theo hướng tiếp cận đã chọn, và (iv) Đóng góp chính của ĐATN là gì, kết quả đạt được sau cùng là gì. Sinh viên cần viết thành đoạn văn, không được viết ý hoặc gạch đầu dòng.
\begin{flushright}
\begin{tabular}{@{}c@{}}
Student\\
\textit{(Signature and full name)}
\end{tabular}
\end{flushright}

\end{document}